% Hodina 49, 1. ročník SPŠ. Rozsah: Rozpis_vyuky_MAT_1_rocnik.tex, řádek 49.

% Plán: 1_rocnik/tematicke_plany/SPS.pdf; společná projektová šablona.

% Učebnice: ucebnice/01_Zakladni_poznatky/Algebraicke_vyrazy/Algebraické výrazy.tex

% Původ úloh: Vlastní zadání a vlastní formulace výkladu; metodická inspirace zdroji.

% Metodický zdroj: externi_zdroje/realisticky/Mnohocleny/07_Rozklad_mnohoclenu_na_soucin_II_vzorce_.pdf

% URL: http://www.realisticky.cz/ucebnice/01%20Matematika%20S%C5%A0/01%20Z%C3%A1kladn%C3%AD%20poznatky/08%20Mnoho%C4%8Dleny/07%20Rozklad%20mnoho%C4%8Dlen%C5%AF%20na%20sou%C4%8Din%20II%20(vzorce).pdf

\documentclass[11pt,a4paper]{article}

% Spolecny styl pro ucitelske pripravy MAT (Work/MAT)
% Verze 1.1, 2026-10-23
\NeedsTeXFormat{LaTeX2e}
\ProvidesPackage{priprava}[2026/10/23 v1.1 Sablona ucitelskych priprav MAT]

% Zakladni typografie a cestina
\RequirePackage[T1]{fontenc}
\RequirePackage[utf8]{inputenc}
\RequirePackage[czech]{babel}
\RequirePackage{lmodern}
\RequirePackage{microtype}
\RequirePackage{amsmath,amssymb}
\RequirePackage{xcolor}
\RequirePackage{enumitem}
\RequirePackage[most]{tcolorbox}
\RequirePackage{tabularx,booktabs,array}
\RequirePackage{fancyhdr}
\RequirePackage{needspace}
\RequirePackage{geometry}
\RequirePackage[hidelinks]{hyperref}

\geometry{a4paper,top=17mm,bottom=18mm,left=23mm,right=23mm}
\definecolor{notegray}{gray}{0.39}
\setlength{\parindent}{0pt}
\setlength{\parskip}{2.5pt plus 1pt minus 0.5pt}
\setlist[itemize]{topsep=2pt,itemsep=1pt,parsep=0pt,leftmargin=1.45em}
\setlist[enumerate]{topsep=2pt,itemsep=2pt,parsep=0pt,leftmargin=1.65em}
\setcounter{secnumdepth}{0}
\renewcommand{\arraystretch}{1.12}

% Zahlavi: vlevo rocnik | globalni tema, vpravo cislo: nazev hodiny
% Zapati zustava zcela prazdne.
\pagestyle{fancy}
\fancyhf{}
\renewcommand{\headrulewidth}{0.35pt}
\renewcommand{\footrulewidth}{0pt}
\setlength{\headheight}{14pt}
\newcommand{\nastavPripravu}[4]{%
  \fancyhead[L]{\small #1\enspace\textbar\enspace #2}%
  \fancyhead[R]{\small #3: #4}%
}

% Minimalisticke cernobile boxy. Kazda definice zvlast.
\tcbset{commonlesson/.style={enhanced,breakable=false,colback=white,colframe=black,
  colbacktitle=white,coltitle=black,boxrule=0.45pt,
  arc=0.8mm,outer arc=0.8mm,left=3mm,right=3mm,top=1.6mm,bottom=1.6mm,
  toptitle=1mm,bottomtitle=0.7mm,fonttitle=\bfseries,
  attach boxed title to top left={xshift=2.5mm,yshift=-1.15mm},
  boxed title style={colback=white,colframe=white,boxrule=0pt,left=1mm,right=1mm}}}
\newtcolorbox{zadaniBox}[1][]{commonlesson,title={Zadání},#1}
\newtcolorbox{prikladBox}[1][]{commonlesson,title={Příklad},#1}
\newtcolorbox{definiceBox}[1][]{commonlesson,title={Definice},#1}

% Spolecny box pro zadani a jeho reseni, bez mezery mezi castmi.
% Pouziti:
%   \begin{prikladReseniBox}
%     Zadani ulohy.
%     \tcblower
%     \begin{align*} ... \end{align*}
%   \end{prikladReseniBox}
% Spodni cast ma jemny sedy podklad a delici linku; lze v ni sazet
% vice radku matematiky. Obe casti maji jen jeden spolecny ramecek.
\tcbset{solutionlesson/.style={
  commonlesson,
  enhanced jigsaw,
  breakable,
  colbacklower=black!3,
  segmentation style={solid,draw=black!25,line width=0.35pt},
  fontlower=\small\color{notegray},
  before lower={\noindent\textbf{Řešení:}\enspace},
  leftlower=3mm,rightlower=3mm
}}
\newtcolorbox{prikladReseniBox}[1][]{solutionlesson,title={Příklad},#1}
\newtcolorbox{zadaniReseniBox}[1][]{solutionlesson,title={Zadání},#1}

% Pomucky pro autora pripravy: sedive poznamky; \reseni pro starsi pripravy
\newcommand{\poznamka}[1]{\par\begingroup\small\color{notegray}\itshape
  \textnormal{\textbf{Poznámka pro učitele.}}\enspace #1\par\endgroup}
\newcommand{\reseni}[1]{\par\begingroup\small\color{notegray}
  \textbf{Očekávané řešení:}\enspace #1\par\endgroup}
\newcommand{\blok}[1]{\needspace{5\baselineskip}\subsection*{#1}}
\newcommand{\vyrok}[1]{\uv{#1}}


\nastavPripravu{1. ročník}{Algebraické výrazy}{49}{Druhé mocniny a rozdíl čtverců}

\begin{document}

\blok{Výklad a souvislosti}

Při roznásobení $(a+b)(a+b)$ vzniknou dva stejné smíšené součiny $ab$. V součinu $(a-b)(a+b)$ se naopak vyruší. Vzorce tak odvodíme z již známého násobení, nemusíme je přijímat bez vysvětlení.

\begin{definiceBox}[title={Definice},breakable=false]
\textbf{Druhá mocnina dvojčlenu} $(a+b)^2$ znamená součin dvou stejných dvojčlenů $(a+b)(a+b)$. Nejde o součet druhých mocnin členů.
\end{definiceBox}


Pro reálná $a,b$ platí
\begin{align*}
(a+b)^2&=a^2+2ab+b^2,\\
(a-b)^2&=a^2-2ab+b^2,\\
a^2-b^2&=(a-b)(a+b).
\end{align*}
První vzorec odvoď roznásobením; ve druhém nahraď $b$ číslem $-b$. Třetí ověř dvěma prostředními součiny s opačnými znaménky. Všechny tři vztahy používáme oběma směry.

\blok{Řešené příklady}

\begin{prikladReseniBox}[title={Příklad 1},breakable=false]
Roznásob $(2x-3)^2$.
\tcblower
Zde $a=2x$, $b=3$. Dostaneme $4x^2-12x+9$. Prostřední člen je $-2\cdot2x\cdot3$, nikoli $-6x$.
\end{prikladReseniBox}


\begin{prikladReseniBox}[title={Příklad 2},breakable=false]
Rozlož $9x^2-25$ a $x^2+6x+9$.
\tcblower
$9x^2-25=(3x-5)(3x+5)$. Druhý výraz je $(x+3)^2=(x+3)(x+3)$, protože prostřední člen je $2\cdot x\cdot3$.
\end{prikladReseniBox}


\blok{Samostatné procvičení}

\begin{zadaniReseniBox}[title={Procvičení},breakable=false]
\begin{enumerate}[label=\alph*)]
\item Roznásob $(a+4)^2$.
\item Rozlož $4y^2-1$.
\item Rozlož $x^2-10x+25$.
\end{enumerate}
\tcblower
\begin{enumerate}[label=\alph*)]
\item $a^2+8a+16$.
\item $(2y-1)(2y+1)$.
\item $(x-5)^2$. Znaménko prostředního členu určuje minus v dvojčlenu.
\end{enumerate}
\end{zadaniReseniBox}


\poznamka{Po každém vzorci nech žáky pracovat samostatně. Zdůrazni, že $a^2+b^2$ se tímto vzorcem nerozkládá. U každého složeného základu vyznač, co zastupuje celé $a$ a celé $b$; jinak se ztrácejí číselné koeficienty.}

\end{document}
